\documentclass[12pt,a4paper]{article}
\usepackage[margin=1in]{geometry}
\usepackage{graphicx}
\usepackage{hyperref}
\usepackage{xcolor}
\usepackage{booktabs}
\usepackage{tabularx}
\usepackage{enumitem}
\usepackage{fancyhdr}
\usepackage{titlesec}
\usepackage{tcolorbox}
\usepackage{fontawesome5}

\definecolor{primary}{HTML}{0A66C2}
\definecolor{accent}{HTML}{057642}
\definecolor{warning}{HTML}{B24020}
\definecolor{bglight}{HTML}{F4F2EE}
\definecolor{darktext}{HTML}{1D2226}

\hypersetup{colorlinks=true,linkcolor=primary,urlcolor=primary}
\pagestyle{fancy}
\fancyhf{}
\fancyhead[L]{\textcolor{primary}{\textbf{LinkedIn Comment Generator}}}
\fancyhead[R]{\textcolor{gray}{User Guide v2.0}}
\fancyfoot[C]{\thepage}

\titleformat{\section}{\Large\bfseries\color{primary}}{}{0em}{}
\titleformat{\subsection}{\large\bfseries\color{darktext}}{}{0em}{}

\newtcolorbox{infobox}{colback=blue!5,colframe=primary,title=\faInfoCircle\ Note,fonttitle=\bfseries}
\newtcolorbox{warnbox}{colback=red!5,colframe=warning,title=\faExclamationTriangle\ Important,fonttitle=\bfseries}
\newtcolorbox{tipbox}{colback=green!5,colframe=accent,title=\faLightbulb\ Tip,fonttitle=\bfseries}

\begin{document}

% ============================================
% TITLE PAGE
% ============================================
\begin{titlepage}
\centering
\vspace*{2cm}
{\Huge\bfseries\color{primary} LinkedIn Comment Generator\par}
\vspace{0.5cm}
{\Large\color{darktext} User Guide \& Operations Manual\par}
\vspace{1cm}
{\large Version 2.0 --- FastAPI + LangGraph Edition\par}
\vspace{0.5cm}
{\large April 2026\par}
\vspace{2cm}

\begin{tcolorbox}[colback=bglight,colframe=primary,width=0.8\textwidth]
\centering
\textbf{AI-powered tool that generates human-like LinkedIn comments}\\[6pt]
CSV Upload $\rightarrow$ Scrape Posts $\rightarrow$ AI Analysis $\rightarrow$ Generate 3 Comments\\[6pt]
\textit{Context-aware $\cdot$ Tone-matched $\cdot$ 44 Comment Angles $\cdot$ RAG-enhanced}
\end{tcolorbox}

\vfill
{\small Built with: Next.js $\cdot$ FastAPI $\cdot$ LangGraph $\cdot$ Groq AI $\cdot$ RapidAPI $\cdot$ SerpAPI\par}
\end{titlepage}

\tableofcontents
\newpage

% ============================================
% GETTING STARTED
% ============================================
\section{Getting Started}

\subsection{What This Tool Does}
The LinkedIn Comment Generator takes a CSV file containing LinkedIn post URLs, scrapes each post's content, analyzes it with AI, and generates 3 human-like comments per post. Each comment uses a different ``angle'' (e.g., sharing experience, asking questions, offering insights) and is tuned to match the post's tone.

\subsection{Access the App}

\begin{itemize}[leftmargin=1.5em]
    \item \textbf{Production:} \url{https://linkedin-comment-gen-production.up.railway.app}
    \item \textbf{Local Development:}
    \begin{itemize}
        \item Frontend: \texttt{http://localhost:3001}
        \item Backend API: \texttt{http://localhost:8000}
    \end{itemize}
\end{itemize}

\subsection{Prerequisites}
\begin{itemize}[leftmargin=1.5em]
    \item A CSV file with LinkedIn post URLs (one per row)
    \item API keys configured (Groq, RapidAPI, SerpAPI)
    \item Commenter profile set up in Settings (for best results)
\end{itemize}

% ============================================
% WORKFLOW
% ============================================
\section{Step-by-Step Workflow}

\subsection{Step 1: Upload CSV}
\begin{enumerate}[leftmargin=1.5em]
    \item Go to the home page (\texttt{/})
    \item Drag \& drop your CSV file or click ``Browse Files''
    \item Supported formats:
    \begin{itemize}
        \item Full LinkedIn post URLs: \texttt{https://linkedin.com/posts/username\_slug...}
        \item Profile URLs: \texttt{https://linkedin.com/in/username}
        \item Bare usernames (single-column CSV): \texttt{justinwelsh}
    \end{itemize}
    \item Click ``Upload \& Extract URLs''
    \item Preview shows detected profiles with POST/PROFILE badges
\end{enumerate}

\begin{tipbox}
Post URLs are preferred over profile URLs. With a post URL, the tool fetches that exact post. With a profile URL, it fetches the latest post (which may not be what you want).
\end{tipbox}

\subsection{Step 2: Scrape Posts}
\begin{enumerate}[leftmargin=1.5em]
    \item Select/deselect profiles using checkboxes
    \item Click ``Scrape Posts''
    \item Tool fetches post data via RapidAPI: author, text, media, likes, comments, shares
    \item \textbf{Cache:} Already-scraped posts are skipped (no duplicate API calls)
    \item Redirects to Dashboard when done
\end{enumerate}

\begin{infobox}
Rate limiting: 1.5 seconds between API calls to respect RapidAPI limits. 10 posts take ~15 seconds.
\end{infobox}

\subsection{Step 3: Dashboard (LinkedIn Feed)}
\begin{enumerate}[leftmargin=1.5em]
    \item Dashboard shows posts in a LinkedIn-style feed layout
    \item Each post card shows: author photo, name, headline, post text, media, reactions
    \item Filter by status: All / Scraped / Analyzed / Generated
    \item Click ``Generate AI Comments'' on any post
\end{enumerate}

\subsection{Step 4: Generate Comments}
When you click ``Generate AI Comments'', the LangGraph pipeline runs:

\begin{enumerate}[leftmargin=1.5em]
    \item \textbf{Analyze} --- Groq AI detects post type, tone, topic, key details
    \item \textbf{Web Search} --- SerpAPI fetches related articles (if post is tactical/controversial)
    \item \textbf{RAG Retrieval} --- Finds 5 high-engagement reference comments from database
    \item \textbf{Generate} --- 3 comments generated \textbf{in parallel} with different angles
    \item \textbf{Humanize + Validate} --- Removes AI phrases, adds natural imperfections, checks quality
    \item \textbf{Retry} --- If quality $<$ 50, regenerates with different angle (max 2 retries)
\end{enumerate}

\subsection{Step 5: View \& Copy Comments}
\begin{enumerate}[leftmargin=1.5em]
    \item Click ``View Generated Comments'' or navigate to \texttt{/post/\{id\}}
    \item See 3 generated comments, each with:
    \begin{itemize}
        \item Angle badge (e.g., \textit{whispered truth}, \textit{hidden cost})
        \item Confidence score (\%)
        \item Word count
        \item Quality bar
    \end{itemize}
    \item Click \textbf{Copy} to copy to clipboard
    \item Use \textbf{thumbs up/down} to rate (improves future generation via feedback loop)
    \item Click \textbf{Regenerate All} for fresh comments
\end{enumerate}

% ============================================
% SETTINGS
% ============================================
\section{Settings --- Commenter Profile}

Navigate to \texttt{/settings} to configure your commenter profile. This is \textbf{critical for quality} --- without it, comments sound generic.

\begin{center}
\begin{tabularx}{\textwidth}{lX}
\toprule
\textbf{Field} & \textbf{Description} \\
\midrule
Name & Your full name (appears in prompt as ``who you are'') \\
Username & LinkedIn username \\
Headline & Your LinkedIn headline (e.g., ``SDE @ Company'') \\
Expertise & Comma-separated skills (e.g., ``AI/ML, product management'') \\
Tone & Writing style: Professional, Casual, Witty, Analytical, Empathetic \\
Avg Comment Length & Target word count (15--100) \\
Common Phrases & Phrases you naturally use, one per line \\
Real Comment Examples & Your actual LinkedIn comments, one per line \\
\bottomrule
\end{tabularx}
\end{center}

\begin{warnbox}
Without real comment examples, the AI has no voice reference. Add at least 5 of your actual LinkedIn comments for best results.
\end{warnbox}

% ============================================
% COMMENT ANGLES
% ============================================
\section{Comment Angles Explained}

The tool uses 21 comment angles. Each produces a different type of comment:

\begin{center}
\small
\begin{tabularx}{\textwidth}{lXl}
\toprule
\textbf{Angle} & \textbf{What It Does} & \textbf{Best For} \\
\midrule
WHISPERED\_TRUTH & Say the quiet part out loud & Tactical posts \\
HIDDEN\_COST & Point out what they're not seeing & Tactical posts \\
SECRET\_NOBODY\_ASKED & Share insider knowledge & Tactical posts \\
LIVED\_IT\_DEEPER & Share harder version of experience & Lessons, stories \\
ME\_TOO\_DEEPER & Specific parallel experience & Lessons, stories \\
MINDSET\_SHIFT & Reframe the situation entirely & Any analytical post \\
INEVITABLE\_NOT\_SURPRISING & Make success feel inevitable & Celebrations \\
CALL\_OUT\_VALUE & Name the specific value & Celebrations, milestones \\
ASK\_QUESTION & Ask genuine tactical question & Engagement posts \\
SHARE\_EXPERIENCE & Parallel experience with data & Any post \\
ADD\_DATA & Contribute facts or metrics & Data-heavy posts \\
EARNED\_INSIGHT & Wisdom from experience & Reflective posts \\
REAL\_RISK & Identify the actual risk & Controversial posts \\
INTENSIFY\_SURGICAL & Amplify with precision & Hot takes \\
REFRAME\_SMARTER & Better perspective & Hot takes \\
VALIDATION & Simple acknowledgment & Any post (short) \\
\bottomrule
\end{tabularx}
\end{center}

% ============================================
% TONE MATCHING
% ============================================
\section{Tone Matching}

The AI detects 7 post writing tones and adapts comment style accordingly:

\begin{center}
\begin{tabularx}{\textwidth}{lX}
\toprule
\textbf{Post Tone} & \textbf{Comment Adapts To} \\
\midrule
Casual/Conversational & Short sentences, informal, no corporate speak \\
Formal/Professional & Polished but still human \\
Listicle/Breakdown & Flowing text (not a list), go deeper on one item \\
Storytelling/Narrative & Respond to emotional core, share parallel story \\
Data-Heavy & Engage with numbers, add context or related data \\
Motivational & Add substance --- specific example or counterpoint \\
Humorous/Witty & Match the energy, be playful \\
\bottomrule
\end{tabularx}
\end{center}

% ============================================
% RAG
% ============================================
\section{RAG --- Reference Comments}

The tool uses 1,719 high-engagement reference comments to improve quality.

\textbf{How it works:}
\begin{enumerate}[leftmargin=1.5em]
    \item When generating comments for an ``advice\_tactical'' post, the system queries the database
    \item Finds top 5 comments that got high engagement on similar post types
    \item Injects them into the prompt as examples
    \item LLM learns from these real examples and produces better output
\end{enumerate}

\textbf{Import RAG data:}
\begin{verbatim}
POST /api/import-rag
Content-Type: multipart/form-data
file: rag_data.csv
\end{verbatim}

\textbf{Check RAG status:}
\begin{verbatim}
GET /api/rag-status
→ {"total": 1719, "by_type": {"advice_tactical": 419, ...}}
\end{verbatim}

\begin{warnbox}
Railway has ephemeral storage. RAG data resets on every deploy. Re-import after each deployment.
\end{warnbox}

% ============================================
% LOCAL DEVELOPMENT
% ============================================
\section{Local Development}

\subsection{Running Locally}

\begin{verbatim}
# Terminal 1: FastAPI backend
cd ~/linkedin-comment-generator/backend
source venv/bin/activate
uvicorn main:app --port 8000 --reload

# Terminal 2: Next.js frontend
cd ~/linkedin-comment-generator
npm run dev -- -p 3001
\end{verbatim}

\subsection{Environment Variables}

Create \texttt{backend/.env}:
\begin{verbatim}
GROQ_API_KEY=gsk_...
RAPIDAPI_KEY=...
RAPIDAPI_HOST=linkedin-scraper-api-real-time-fast-affordable
              .p.rapidapi.com
SERPAPI_KEY=...
\end{verbatim}

\subsection{Viewing Logs}
All pipeline steps are logged in the FastAPI terminal with colored output:
\begin{itemize}[leftmargin=1.5em]
    \item \texttt{[Analyzer]} --- Post analysis results
    \item \texttt{[Groq]} --- Full prompts and responses with token counts
    \item \texttt{[RAG]} --- Reference comment retrieval
    \item \texttt{[Humanizer]} --- Changes applied
    \item \texttt{[Generator]} --- Angle selection, validation, confidence breakdown
\end{itemize}

\newpage
% ============================================
% COSTING
% ============================================
\section{Monthly Costing}

\subsection{Per-Post Cost Breakdown}

\begin{center}
\begin{tabularx}{\textwidth}{Xlrr}
\toprule
\textbf{Service} & \textbf{Calls/Post} & \textbf{Cost/Call} & \textbf{Cost/Post} \\
\midrule
RapidAPI LinkedIn Scraper & 1 & \$0.025 & \$0.025 \\
Groq llama-3.3-70b (analysis) & 1 & \$0.0014 & \$0.0014 \\
Groq llama-3.3-70b (generation $\times$3) & 3 & \$0.0014 & \$0.0042 \\
Groq llama-4-scout (vision) & 0.5 avg & \$0.0004 & \$0.0002 \\
SerpAPI (web search) & 0.5 avg & \$0.015 & \$0.0075 \\
\midrule
\textbf{Total per post} & & & \textbf{\$0.038} \\
\bottomrule
\end{tabularx}
\end{center}

\begin{infobox}
Groq AI is extremely cheap (\$0.59/M input tokens). The biggest cost is RapidAPI and SerpAPI. Vision calls are almost free (\$0.11/M tokens).
\end{infobox}

\subsection{Monthly Cost by Volume}

\begin{center}
\begin{tabular}{rrrrrr}
\toprule
\textbf{Posts/Mo} & \textbf{API Cost} & \textbf{RapidAPI Plan} & \textbf{Hosting} & \textbf{SerpAPI Plan} & \textbf{Total/Mo} \\
\midrule
100 & \$3.80 & Free tier & \$5 & \$25 (starter) & \textbf{\$33.80} \\
500 & \$19.00 & \$25 (Pro) & \$5 & \$25 (starter) & \textbf{\$74.00} \\
1,000 & \$38.00 & \$25 (Pro) & \$5 & \$75 (developer) & \textbf{\$143.00} \\
5,000 & \$190.00 & \$75 (Ultra) & \$20 & \$150 (production) & \textbf{\$435.00} \\
\bottomrule
\end{tabular}
\end{center}

\subsection{Cost Optimization Tips}
\begin{itemize}[leftmargin=1.5em]
    \item \textbf{Cache aggressively} --- already-scraped posts skip API calls (saves RapidAPI costs)
    \item \textbf{Skip web search} for celebration/hiring posts (saves SerpAPI costs)
    \item \textbf{Groq is cheap} --- don't optimize here, it's \$0.006/post for 4 LLM calls
    \item \textbf{Batch processing} --- upload large CSVs to amortize per-session overhead
\end{itemize}

\end{document}
